% resumo em português
\setlength{\absparsep}{18pt} % ajusta o espaçamento dos parágrafos do resumo
\begin{resumo}
Os Saberes e Sistemas Agrícolas Tradicionais no Território de Identidade Semiárido Nordeste II (Bahia) enfrentam falhas de mercado decorrentes de assimetria informacional, onde a ausência de mecanismos formais de mensuração e sinalização impede a captura de valor premium associado a ativos intangíveis englobando autenticidade, sustentabilidade e reputação territorial. Essa lacuna desdobra-se em fricções institucionais e custos de transação que criam barreiras de entrada para comunidades com baixa capacidade técnico-institucional, sustentando subprecificação e vulnerabilidade à apropriação indevida. Este estudo desenvolve e valida um sistema integrado de governança bioeconômica, suportado por modelagem computacional, visando a mensuração de ativos intangíveis e a mitigação dessas assimetrias. A pesquisa estrutura-se em cinco módulos articulados, abrangendo diagnóstico institucional e regulatório para mapear gargalos de governança e custos de transação, elicitação estruturada de variáveis mediante método Delphi convertendo saberes tácitos em critérios mensuráveis e consensuados, modelagem fuzzy baseada na Teoria dos Conjuntos Difusos para construção do Índice de Valoração Bioeconômica capturando ambiguidade inerente a dimensões culturais e simbólicas, validação de mercado através de experimentos de escolha e testes de campo para mensurar disposição a pagar e testar o índice como mecanismo de signaling, e design participativo incorporando usuários finais na co-criação de regras e interfaces para garantir apropriação e adoção continuada. O sistema proposto integra inteligência territorial mediante sensoriamento remoto e geoprocessamento, governança de dados por protocolos de consenso e rastreabilidade, e inferência fuzzy produzindo métricas auditáveis de apoio à decisão em cenários de acesso a mercado, posicionamento estratégico e proteção de propriedade intelectual. O resultado esperado consiste na consolidação de um ecossistema de pesquisa aplicada com capilaridade formativa, capaz de orientar trilhas empíricas e computacionais interdisciplinares e gerar soluções tecnológicas territorialmente ancoradas e replicáveis para contextos de bioeconomia comunitária.

\textbf{Palavras-chave}: Governança Bioeconômica; Lógica Fuzzy; Assimetria Informacional; Ativos Intangíveis; Saberes Tradicionais.
\end{resumo}

